\section{Discussion}\label{sec:discussion}

\para{Beyond binding}
We expect the no-drop idea to transfer to any LLM auditor that keeps
intermediate notes: a withdrawn or relocated suspicion deserves at least as
much scrutiny as the suspicion itself, and structured intermediate output is
what makes ``never seen'' and ``seen and lost'' distinguishable when the
agent is evaluated.

\para{Lessons for developers}
The signature analogy doubles as a checklist. \emph{Sign what you act on}:
every recipient, amount and fee is a public input or hashed into one that the
consumer recomputes (P1). \emph{Separate domains}: commit to chain, contract
and policy, and check roots and prior states against the consumer's own
storage (P2). \emph{Derive nonces from the statement} and reduce them into the
field before storing (P3). \emph{Pin the key} and check the verifier's result
(P4). \emph{Test with a malicious caller}: an honest proof submitted by a
second account with changed arguments must fail. This single test, the
invariant of Section~\ref{sec:design-loop}, is what turned the fuzzer from 0
to 5 of 13 flaws.

\para{Toward binding by construction}
Binding flaws exist because the public-input vector is an untyped array whose
meaning is restated by hand in the circuit, the verifier wrapper and the
consumer. Circuit languages could export a \emph{typed statement} schema from
which verifier generators emit a typed consumer interface, removing P5 and
making P1 visible. A standard \emph{effect commitment} over recipient,
amount, chain, contract and nonce, computed in the circuit and recomputed in
the consumer, would play the role EIP-712~\cite{eip712} plays for signatures
and address P1--P3 together. Verifier libraries could make the key or program
ID an immutable constructor argument (P4). \DEMO{Applied to the 87 core
findings, typed statements would have prevented 41\% and effect commitments
63\%.} The binding question---does the verified object commit to every value
the effect uses?---also applies beyond ZK: to signed permits that omit the
spender, TEE attestations whose report data omits the request, and bridge
messages that omit the destination domain.

\subsection{Threats to validity}\label{sec:threats}

\para{Selection and labels}
We study public audits of projects that chose to publish them; unaudited
projects may differ. Labels were proposed by an LLM and adjudicated by two
raters who discussed every item, so they are not independent annotations; a
written guide, a code-level pass for every evaluated item \DEMO{and a blind
relabel ($\kappa=0.78$)} mitigate this, and we release every label with its
notes. Auditors' locations are one of several valid ones; mechanism-level
recall bounds the effect.

\para{Development bias}
The post-stage rules were written while studying dev misses, and one
admissibility rule and one catalogue fact were informed by an item in the
original check split, which is among the gains; dev recall is therefore
optimistic, and we report it only as development evidence. The keep rule was
fixed after a preliminary single-run measurement of no-drop and
multi-location and before any other step was measured; one counting detail
was corrected afterwards, which turned the knowledge step from pass to fail
and changed no other decision (Appendix~\ref{app:count}). Held-out claims
rest on sets the building process never saw.

\para{Location credit and judge}
Multi-location reporting raises the chance of matching the auditor's
location. The judge's mechanism check, 0 of 17 fixed twins flagged at
strict, and per-finding counting bound this effect but do not remove it. The
judge is an LLM with a frozen prompt and about 5\% verdict instability; all
configurations share cached verdicts.

\para{Variance and contamination}
With 28 dev flaws and two runs, intervals are wide and 8 of 18 hits fragile;
held-out numbers are medians of three runs. Audit reports may be in the
model's training data. Mutants are new code, snapshots contain no audit text,
and fixed twins are a partial control: a model recalling a report would also
flag the patched code, and v4 flags none at strict. \DEMO{On round~3, whose
reports were all published after the model's training cutoff, v4 finds 3 of 7,
as on the check set, and the three new flaws were unknown before our scans.}

\para{Leakage}
The coordinating author saw v2's test results, so B1 is development data for
v4. Every version was built by a separate process restricted to development
data, test benchmarks live in a separate repository, and every exposure is
logged (Appendix~\ref{app:timeline}).

\subsection{Limitations}\label{sec:limits}
\tool{} needs an LLM and costs about one to two dollars per scan, and all
runs use one model. Its helpers are syntactic, so library code that is not
vendored is a blind spot, and the residue of RQ3 sits exactly there. The
loop targets EVM consumers: flaws whose effect lies in a zkVM host or off
chain stay at E1, and real-proof tests need proving artifacts that most
projects do not ship. Binding is relative to the intended policy $\Pi$, which
\tool{} infers from code and documentation and can get wrong.
